
\(P_A(x) = \text{det}(xI_n - A) \in R[x]\) it holds \(P_A(A) = 0\).

Applications
With this we can d

Define an \(R\)-module \(M\)

The axiomatic definition is the existence of two binary operations

\(+ : M \times M \to M, (x, y) = x + y\)
\(\cdot : R \times M \to M, (r, x) = rx\)
satisfying
\((M, +)\) is an abelian group.
and for \((a, b \in R, x, y \in M)\) it holds
\(a(x + y) = ax + ay\)
\((a + b)x = ax + bx\)
\((ab)x = a(bx)\)
\(1x\)
for a (left) \(R\)-module \(M\), 
In case of a right \(R\)-module, \(a, b \in R\) have to be to the right
of the \(x, y \in M\) in the axioms.

Alternatively a module can be seen as a ring homomorphism \(\phi: R \to \text{End}(M)\)
that is \(ax = \phi(a)(x)\), which naturally satisfyies the axioms above. 
One calls \(\phi\) a ring action on \(M\).

A left \(R\)-module can be turned into a right \(R^{\text{op}}\)-module by means of the anti-isomorphism
(it's actually not an isomorphism, just a bijective map) \(\text{op} : R \to R^{\text{op}}\) \(\text{op}(ab) = ba\),
by the map \(\psi = \phi\circ\text{op} : R \to \text{End}(M)\), which actually induces a morphism \(\tilde{\psi} : R^{\text{op}} \to \text{End}(M)\).

\iffalse
by the induced morphism \(\phi^{\text{op}} : \), \(\phi^{\text{op}} = \psi\text{op_M}\phi\text{op}^{-1}\),
\text{op_M}
    where


\(\text{op}(\phi)(a)(x) = x\phi(a)\)
\fi


What is the natural relation between abelian groups
and modules?


Take any abelian group \(G\) and make it to a \(\mathbb{Z}\)-module
by \(\phi : \mathbb{Z} \to \text{End}(G)\), \(\phi(z)(g) = \sum_{i = 1}^{|z|} \text{sgn}(z)g\), in fact this
is the only possible \(\mathbb{Z}\)-module structure on \(G\).


Define an \(R\)-module homomorphism or \(R\)-linear map

if the map \(\phi : M \to N\) satisfies for all \(x, y \in M\) and \(a \in R\)
\(\phi(x + y) = \phi(x) + \phi(y)\)
\(\phi(ax) = a\phi(x)\)

then \(\phi\) is an (left) \(R\)-module homomorphism. 

The second property is sometimes tricky e.g. \(\phi : \mathbb{R}^2 \to \mathbb{R}^2, \phi(x, y) = (x, 2y + 1)\) is not a homomorphism
as it fails to respect scalar multiplication in the second component.

If \(R\) is commutative then \(\text{Hom}_{R}(M, N)\) is an \(R\)-module.


Define a submodule \(N\) of and \(R\)-module \(M\)

Axiomatic definition

\(N \neq \emptyset\)
\(x, y \in N \implies x + y \in N\)
This can be summarized with \(N\) being a subgroup of \(M\)
\(x \in N \implies ax \in N, \forall a \in R\)

Then \(N\) is itself an \(R\)-module and the inclusion \(\iota : N \to M\) 
is an \(R\)-module hom.
The quotient \(R\)-module \(M/N\) is defined as usual via, \(\alpha\pi(x) = \alpha x + N\).

State the isomorphism theorems for modules
Let \(\phi: M \to N\) a \(R\)-module hom.
then \(M/\text{ker}(\phi) \simeq \text{im}\phi\)
(In particular \(\text{ker}\phi, \text{im}\phi\) are both submodules.

For \(N \subset M \subset L\) it holds
\( (L/N)/(M/N) \simeq L/M\)
For \(N_1, N_2 \subset M\) 
\((N_1 + N_2)/N_1 \simeq N_2/(N_1 \cap N_2)\)


Define the direct sum and product of modules
and state their properties


Let \(\{M_i\}_{i \in I}\) be a family of 
\(R\)-modules. Define the direct product by
\(\Pi_{i \in I} M_i = \{(x_i)_{i \in I} | x_i \in M_i\}\)
The direct product inherits an 
\(R\)-module structure by defining the operations coordinatewise.
Equivalently elements in \(\Pi_{i \in I} M_i\) are maps \(f : I \to \bigcup_{i \in I} M_i\)
s.t. \(f(i) \in M_i, \forall i\).

Direct sums of modules 
\(\bigoplus_{i \in I} M_i\)
are defined in the same way except only finitely many of the coordinates of the tuples are nonzero. 
Thus if \(I\) is finite direct sums and products are the same.
We say \(M\) is the internal direct sum of the family if 
\(M = \sum M_i\) and \(M_j \cap \sum_{i \neq j} M_i = \{0\}\)

They satisfy the usual universal property 

Direct product
For an \(R\)-module \(N\) and \(\phi_i : N \to M_i\) for each \(i\),
there is unique hom. \(\theta : N \to \Pi_{i in I} M_i\) which factors through
\(\pi_i \circ \theta(n) = \phi_i(n)\) for the canonical projection \(\pi : \Pi_{i in I} M_i \to M_i\).

Direct sum
For an \(R\)-module \(N\) and \(\phi_i : M_i to N\) for each \(i\),
there is unique hom. \(\theta : \bigoplus{i in I} M_i \to N\) which factors through
\(\theta \circ \iota (m) = \phi_i(m)\) for the canonical embedding \(\iota : M_i \to \bigoplus{i in I} M_i\).


Define the ideal quotient of submodules \(I, J \subset M\) and the annulator. State also some essential properties

Definition
\((I : J) = {a \in R | aJ \subset I } \subset R\)
\((0 : M) = \text{Ann}_R(M)\)
Properties
\((I : J) = \text{Ann}_R(J + I)/I\)
\((I : JK) = ((I : J) : K)\) 
\((I : R) = I\) 
\((I : J + K) = (I : J) \cap (I : K)\) 
\((I : (r)) = \frac{1}{r}(I \cap (r))\) if \(R\) is an integral domain.

Let \(I \subset \text{Ann}_R(M)\) an ideal, then \(M\) is also a \(R/I\)-module
further \(\text{Ann}_{R/\text{Ann}_R(M)}(M) = 0\).

If \(R\) is commutative then 
(Sometimes injectivity is refferd


Define a generating set, linear independence and free modules

For an \(R\)-module \(M\) and a family \(\{x_i\}_{i \in I} \subset M\)
is called a generating set if for each \(m \in M\)
\(m = \sum_{i \in I} \alpha_i x_i\) with only finitely many \(\alpha_i \neq 0\).

If \(|I| < \infty\) we say that \(M\) is finitely generated.

It is lin. if 
\(\sum{i \in I} \alpha_i x_i = 0 \iff \alpha_i = 0, \forall i \in I\)

If \(\{x_i\}_{i \in I}\) satisfies both these conditions its a basis 
and \(M\) is free \(R\)-module.

Equivalently, if there is module homomorphism \(\phi R^{I} \to M\) with \(\phi(f) = \sum_{i \in I} f(i) x_i\)
which is surjective (gen. set) and injective (lin. ind.) 

Equivalently this is a homomorphism from 

How are ideals in \(I \subset R\) related to quotient modules of \(M\)?

If \(I\) is a (left) ideal in \(R\), then \(M/IM\) is an \(R/I\)-module.
Moreover a basis of \(M\) \(\{x_i\}_{i \in I}\), is a basis for \(M/IM\) via \(\{\pi(x_i)\}_{i \in I}\).

An important application is the following:
Let \(R\) be a non trivial commutative ring and \(M\) a fin. gen. free \(R\)-module, then for any two bases
\(\{x_i\}_{i \in I}\), \(\{x_j\}_{j \in J}\), \(|I| = |J|\).

Proof:
Every commutative ring has at least one maximal ideal, let it be \(I\), then \(M/IM\) is an \(R/I\)-vector space (as \(R/I\) is a field)
with bases \(\{\pi(x_i)\}_{i \in I}\), \(\{\pi(x_j)\}_{j \in J}\), but any 2 bases of a vector space must have the same cardinality.


State and proof the Cayley Hamilton theorem, and describe some of its uses.

Let \(R\) be a commutative ring \(A = (a_{i,j}) \in R\) a \(n \times n\) matrix, for the characteristic polynomial
\(P_A(x) = \text{det}(xI_n - A) \in R[x]\) it holds \(P_A(A) = 0\).

\(A = \phi : R^n \to R^n\), that is we can represent \(A\) is a module homomorphism, further \(R[\phi] \subset \text{Mat}_n(R)\) generated
by the identity \(I_n\) and \(A\). Let \(\{e_i = (0, \dots, 1_i, 0 \dots) \}\) be the canonical basis, then we have the linear system of eqs.
\(\phi(e_i) = \sum^{n}_{j = 1} a_{i, j} e_j\) which is equivalent to \(\sum_{j = 1}^{n} (\delta_{i,j} \phi - a_{i, j}) e_j = 0\), for each \(i \in 1 \dots n\) (FC),
define the matrix \(B = (\delta_{i,j} \phi - a_{i, j})_{i,j}\), note that \(P_A(x) \to P_A(A) = \text{det} B\) (by the morphism \(R[x] \to R[\phi]\)).
With the adjugate matrix \(B^{*}B = \text{B}I_n\) we can deduce \(0 = B^{*}B(e_1, e_2, \dots)^{T} = (\text{det}Be_1, \dots)^{T}\) where the first equality stems from FC.

We can find \(A^n\) in terms of lower powers according to this theorem (as \(p(A) = A^n + c_{n-1}A^{n-1} \dots + c_0 I_n\)).
For \(R\) a field a non-trivial equivalence is that the minimal polynomial of \(A\) divides its characteristic polynomial.

What is the determinant trick?

Let \(R\) be a commutative ring, \(M\) a fin. gen. \(R\)-module, \(I \subset R\) and ideal and \(\phi : M \to M \in \text{End}(M)\) s.t. \(\phi(M) \subset IM\) 
then there is some relation

\(\phi^{n} + a_{n-1}\phi^{n-1} + \dots + a_0 = 0\) with \(a_i \in I\).

The proof is essentially the same as for the Caley-Hamilton theorem with the most table difference that we evaluate the determinant to get this expression
and use the fact that \(\phi(M) \subset IM\) so that all entries in the matrix \(B\) are in \(I\).

State and proof the Lemma of Nakayama:

Let \(R\) be a commutative ring, \(M\) a fin. gen. \(R\)-module, \(I \subset R\) an ideal and \(M = IM\) then there is some \(r \in R\) with
\(r \equiv 1 \text{mod} I\), s.t. \(rM = 0\). This follows from the determinant trick with \(\phi = \text{id}_M\) and \(x = 1 + a_{n-1} + \dots + a_0\).

Lemma of Nakayama
Let \(R\) be a local ring with maximal ideal \(\mathfrak{m}\). When \(M\) is fin. gen. with \(M = \mathfrak{m}M\) then \(M = 0\)
By the previous we have an \(x \in R\), \(r \equiv 1 \text{mod} \mathfrak{m}\) and \(x \in \text{Ann}_R(M)\). Suppose \(\text{Ann}_R(M) \neq R\)
then \(\text{Ann}_R(M) \subset \mathfrak{m}\) by maximality and from the conruence \(1 \in \mathfrak{m}\), a contradiction.

Some useful consequences are then found for decompositions of modules

Let \(R\) be a local ring with maximal ideal \(\mathfrak{m}\). When \(M\) is fin. gen. and \(N \subset M\) with \(M = N + \mathfrak{m}M\) then \(N = M\).
Proof: Apply Nakayama's Lemma to \(M/N\) and note that \(\mathfrak{m}(M/N) = (\mathfrak{m}M + N)/N = M/N\).

Also if \(\{\pi(x_i)\}_{i \in I}\) are a basis of the \(R/\mathfrak{m}\)-vector space \(M\), so is \(\{x_i\}_{i \in I}\) a basis of the \(R\)-module \(M\).
Proof: Let \(N = \text{span}(\{x_i\}_{i \in I})\) be the submodule generated by these elements, then \(M = N + \mathfrak{m}M\) from the splitting SES \(0 \to N \to M \to \mathfrak{m}M \to 0\) 
and by the previous theorem the result follows.


Characterize noetherian modules
and give some examples and theorems

Let \(M\) be an \(R\)-module, tfae.

1. Every submodule of \(M\) is finitely generated.
2. Every ascending chain of submodules \(N_1 \subset N_2 \subset \dots\) becomes stationary (\(N_i = N_j, j \geq i\) for some \(i\)), also called the ACC-condition.
3. Every non empty set of submodules of \(N\) has a maximal element.

Examples:
Every PID ring is noetherian. \(C[0, 1]\) is not noetherian 

Properties:
Every submodule of a noetherian module is noetherian.
If \(f : M \to N\) is surjective and \(M\) noetherian, so is \(N\).

For a SES \(0 \to M' \to M \to M'' \to 0\) it holds that \(M\) is noetherian if and only if \(M', M''\) are.
It follows that \(M\) is noetherian if for a noetherian submodule \(N\) also \(M/N\) is noetherian.

Let \(R\) be a left-noetherian ring and a \(M\) fin. gen. module over it. Then \(M\) is also noetherian.

State an important lemma which contraposits noetherian modules.

If \(M, N\) are \(R\)-modules with \(M \simeq M \oplus N\) and \(N \neq 0\), then \(M\) is not noetherian

The proof is based on the idea to take the poset \(\Sigma\) of \(\{ (M', N') | M' \oplus N' \simeq M \}\) and in particular
\(M' \simeq M\), assuming there would be a maximal element \((M', N')\) then \(M \simeq M' \oplus N'\) but also \(M' \simeq M \oplus N\) by assumption
thus \(M \simeq M \oplus N' \oplus N\), since \(N \neq 0\), \(N' \subset N' \oplus N\) (one needs to use the explicite isomorphism to make the notation rigorous)
\(N'\) is not maximal, a contradiction.

This in turns implies that for two different bases \(\{x_i\}_{i \in I}\), \(\{y_j\}_{j \in J}\) of a noetherian module, we must have \(|I| = |J|\),
since w.l.o.g. \(|J| \geq |I|\), \(M \simeq R^{|I|}\), so \(M \simeq R^{|I|} \oplus R^{|J| - |I|}\), but by the previous \(R^{|J| - |I|} = 0\)
as \(M \simeq R^{|I|}\)

